\documentclass[11pt,a4paper]{article}
\usepackage{fullpage}
\usepackage{amsmath}
\usepackage{amssymb}
\usepackage{textcomp}
\usepackage[utf8]{inputenc}
\usepackage[T1]{fontenc}
\textheight=10in
\pagestyle{empty}
\raggedright
\usepackage[left=0.85in,right=0.85in,bottom=0.75in,top=1in]{geometry}
\usepackage{hyperref}

\hypersetup{
  colorlinks=true,
  urlcolor=black,
  linkcolor=black,
  citecolor=black,
  pdfauthor={Valerie L. Baldinger},
  pdftitle={Curriculum Vitae - Valerie L. Baldinger}
}
\urlstyle{same}

% Compact spacing keeps all three complete paper entries on page one.
\linespread{1.08}

\newcommand{\header}[1]{
  {\hspace*{-18pt}\vspace*{6pt}\textsc{#1}}
  \vspace*{-6pt}\lineunder
  \vspace{2mm}
}

\newcommand{\lineunder}{
  \vspace*{-8pt}\\
  \hspace*{-18pt}\hrulefill\\
}

\begin{document}
\vspace*{-52.5pt}

%==== Profile ==== %
\vspace*{-10pt}
\begin{center}
  {\LARGE\scshape Valerie L. Baldinger}\\[2pt]
  \small
  \href{mailto:valerie.baldinger@nyu.edu}{valerie.baldinger@nyu.edu} $\cdot$
  \href{https://valeriebaldinger.github.io/}{valeriebaldinger.github.io} $\cdot$ 646-726-8149\\[2pt]
  NYU Stern $\cdot$ 44 West 4th Street, 9-193F $\cdot$ New York, NY 10012\\[1pt]
  Citizenships: USA, Austria
\end{center}
\vspace{-1mm}

%==== Education ==== %
\header{Education}
\textbf{PhD in Finance, NYU Stern}\hfill 09/2022 - present\\
{\small\textit{Committee: Arpit Gupta (Chair), Toomas Laarits, Thomas Philippon, Indira Puri, Olivier Wang}\par}
\vspace{3mm}

\textbf{MSc in Economics, London School of Economics}\hfill 08/2018 - 06/2020\\
Graduated \textit{with distinction}
\vspace{3mm}

\textbf{BA in Econ., Philosophy \& Pol. Sci., University of Zurich}\hfill 08/2015 - 06/2018\\
\vspace{3mm}

%==== Working Papers ==== %
\header{Working Papers}
``Generative AI and Retail Investment''\\
\textit{Job Market Paper}\\
\vspace{2mm}
{\small
Generative AI is a general purpose technology with rapidly growing adoption among retail investors. How does the adoption of this technology change their investment behavior? I study this question using transaction-level data linking investors' prompts and AI answers to their trades and portfolio holdings. Using a staggered difference-in-differences design, I find that generative AI adoption increases investment in index funds and reduces investment in individual stocks, thereby improving diversification and portfolio efficiency. These improvements are larger for investors who were less diversified before adoption.\par
}
\vspace{3mm}

``LLM Agents as Traders in Experimental Asset Markets''\\
\vspace{2mm}
{\small
A large literature is concerned with understanding why asset price bubbles form. The theoretical literature has proposed multiple mechanisms, including extrapolative beliefs, heterogeneous expectations, and speculation as possible explanations. I develop a methodology using LLM agents as traders in experimental asset markets to study why asset price bubbles form. The methodology exploits unique properties of LLM subjects to disentangle the contribution of mechanisms proposed in the literature: their memories can be selectively edited without awareness, their reasoning can be probed without contaminating subsequent behavior, and market conditions can be varied at scale. I find that LLM traders generate many of the same market phenomena observed with human subjects: price bubbles emerge and crash, experience reduces mispricing, short selling attenuates overpricing, and professional personas trade more efficiently than student personas. These results hold across variations in trading mechanisms and participant characteristics.\par
}
\vspace{3mm}

%==== Work in Progress ==== %
\header{Work in Progress}

``Lowering the Cost of Entry: Evidence from Firm AI Adoption in Procurement Auctions''\\
with Dean Parker\\
\vspace{2mm}
{\small
Generative AI, as a labor-augmenting technology, has the potential to raise productivity across a broad range of tasks, with implications for small firms' competitiveness. In this paper, we study whether the proliferation of generative AI has changed competition for government contracts. Bidding for these contracts involves labor-intensive document analysis and preparation, tasks for which generative AI is well suited and widely used. We assemble a dataset of state-level procurement contracts, bids, and award outcomes across industries and government departments. We find an upward level shift in bidding participation following the emergence of generative AI, particularly for smaller contracts. This is consistent with AI lowering analysis and bid-preparation costs, thus lowering the barrier to entry into procurement auctions. To examine this mechanism, we exploit cross-state variation in exposure to generative AI stemming from differences in the amount of documentation required to submit a bid. We also provide direct evidence on AI use in procurement auctions by using an AI detection tool on firms' bid documents. We examine which firms adopt AI, how AI use varies across contract types, and how it is associated with the likelihood of winning a contract.\par
}

%==== Work Experience ==== %
\newpage
\vspace*{-48pt}
\header{Work Experience}
\textbf{Senior Research Analyst}, Federal Reserve Bank of New York\hfill 07/2021 - 09/2022\\
\vspace{1mm}

\textbf{Junior Research Scientist}, Department of Finance, NYU Stern\hfill 06/2020 - 06/2021\\
\vspace{1mm}

\textbf{Research Assistant}, Department of Economics, London Business School\hfill 05/2019 - 03/2020\\
\vspace{2mm}

%==== Teaching Experience ==== %
\header{Teaching Experience}
\textbf{NYU Stern}\\
\vspace{2mm}
\textit{Instructor}\\
Foundations of Finance, Undergraduate (Summer 2025)\hfill Average instructor rating: 5.00\\
\vspace{2mm}

\textit{Teaching Assistant}\\
{\small
\begin{tabular*}{\linewidth}[t]{@{}l@{\extracolsep{\fill}}l@{}}
Foundations of Finance, Undergraduate (Spring 2025, Fall 2025) & Instructor: Prof. Toomas Laarits\\
Finance Concepts and Math, MBA (Summer 2025) & Instructor: Prof. Alexi Savov\\
Foundations of Finance, MBA (Summer 2025) & Instructor: Prof. Alexi Savov\\
\end{tabular*}\par
}
\vspace{2mm}

%==== Presentations ==== %
\header{Presentations}
NYU Stern Finance Seminar (2024, 2025, 2026)\\
Johns Hopkins Carey Finance Conference (2025)
\vspace{2mm}

%==== Awards ==== %
\header{Awards}
Jules I. Bogen Fellowship, NYU Stern Finance (2026--2027)\\
David M. Graifman Merit Award, NYU Stern Finance (2026)\\
Martin J. Gruber Best 2nd Year Paper Award, NYU Stern Finance (2024)
\vspace{2mm}

%==== Grants ==== %
\header{Grants}
Pangram Research Grant, \$5,000\\
Google AI Research Grant, \$1,000
\vspace{2mm}

%==== Programming Languages ==== %
\header{Programming Languages}
Python, Stata, MATLAB, LaTeX

\end{document}
